% Generated draft: requires mathematical and source review.
\documentclass[11pt]{article}
\usepackage[margin=25mm]{geometry}
\usepackage{fontspec}
\setmainfont{Latin Modern Roman}
\usepackage{amsmath,amssymb,mathtools,unicode-math}
\setmathfont{Latin Modern Math}
\usepackage{enumitem,needspace,xurl,hyperref}
\setlength{\parindent}{0pt}
\setlength{\parskip}{5pt}
\setlength{\emergencystretch}{3em}
\newcommand{\N}{\mathbb{N}}
\newcommand{\Z}{\mathbb{Z}}
\newcommand{\Q}{\mathbb{Q}}
\newcommand{\R}{\mathbb{R}}
\newcommand{\C}{\mathbb{C}}
\newcommand{\F}{\mathbb{F}}
\newcommand{\A}{\mathbb{A}}
\newcommand{\PP}{\mathbb{P}}
\newcommand{\E}{\mathbb{E}}
\newcommand{\eps}{\varepsilon}
\newcommand{\dd}{\,\mathrm{d}}
\begin{document}
\section*{UnsolvedMath ID: 30001014}
\textbf{Source problem code:} catalogue entry
\subsection{Definitions and mathematical statement}
A maximal abelian self-adjoint subalgebra (MASA) $C\subset A$ of a $C^*$-algebra is an abelian $C^*$-subalgebra maximal among abelian $C^*$-subalgebras of $A$. Its Gelfand spectrum $\widehat C$ is the space of nonzero characters with the Gelfand topology. When $C$ is unital, $\widehat C$ is compact Hausdorff and $C\cong C(\widehat C)$.

The source asks whether every compact Hausdorff space $X$ can be realized as $\widehat C$ for a MASA inclusion $C\subset A$ exhibiting the non-maximal-abelian tensor-product pathology of the report's Proposition 4. Thus the intended quantifiers are
\[
\forall X\text{ compact Hausdorff}\quad
\exists(C\subset A)\text{ a MASA inclusion with }\widehat C\cong X
\text{ and the specified Proposition 4 pathology}.
\]
The ambient algebra and all further tensor-product data may depend on $X$; no separability or metrizability restriction is supplied.

For orientation only, a common kind of such pathology uses another MASA $D\subset B$ and a nonminimal $C^*$-tensor norm $\alpha$ such that the image of $C\otimes D$ is abelian but not maximal abelian in $A\otimes_\alpha B$. Non-maximality can be expressed by the existence of a self-adjoint element of the ambient tensor product commuting with this image but lying outside it. The supplied record does not reproduce Proposition 4, so it does not establish that this particular two-inclusion schema, self-tensoring, a fixed second algebra, or a particular nonminimal norm is the intended requirement. They must not be silently substituted for one another.

The precise unitality conditions, admissible tensor products and any restrictions on $X$ in that proposition are essential missing information. In particular, the literal ``every compact Hausdorff space'' includes finite, nonmetrizable and possibly empty spaces under the usual convention; no exclusion or modification is assumed here.

The supplied literature assessment checked 2026-08-21 labels this realization question open, says the report poses it without a solution and reports no later general realization theorem located. This is dated source curation, not an independently verified conclusion.
\subsection{Short English statement}
Ask whether every compact Hausdorff space is the spectrum of a MASA inclusion with the tensor-product failure specified in Proposition 4. The proposition's tensor norm, second algebra and admissibility conditions are missing, so the pathology cannot be uniquely instantiated.
\subsection{Sources}
\begin{itemize}
\item Local UnsolvedMath registry, permanent ID 30001014, source code catalogue entry; supplied complete record, including its Proposition 4 reference and corrected-verified editorial-cleanup note. Literature checked 2026-08-21: the assessment is supplied dated curation, not independently verified here.
\item $C^*$-Algebras (2009), Oberwolfach Reports, DOI 10.4171/owr/2008/37, pp. 2095--2164: \url{https://doi.org/10.4171/owr/2008/37}. The source title's damaged escaping is normalized typographically; proposed-year metadata gives 2008.
\item Wassermann, Tensor products of maximal abelian subalgebras: \url{https://arxiv.org/abs/0711.0131}.
\end{itemize}
\end{document}
